% !PRINTED 6
\weberpage{6}
%<<CONT
and call this a \gesperrt{linear substitution}. If we interpret $x$ and $x'$, as in § 1, as points on two straight lines $L$ and $L'$, then (2) fixes a \gesperrt{mutually one-to-one relation} between the points of $L$ and $L'$, that is, a \gesperrt{mapping}. To the values $x=\infty$ and $x'=\infty$ there correspond the infinitely distant parts of the straight lines, which we likewise regard as definite points. The point $x'=\infty$ then corresponds to the point $x=\alpha:\gamma$, and the point $x=\infty$ to the point $x'=-\delta:\gamma$; and only when $\gamma=0$, that is, when the substitution (2) is \gesperrt{integral}, do the infinitely distant points on $L$ and $L'$ correspond to one another.

The substitution (2) is not altered if the four transformation numbers are multiplied by the same factor. The determinant $r$ is multiplied by the square of this factor, and one can therefore choose this factor so that the determinant takes any prescribed value, for example the value 1.

If we apply the substitution (2) to any two points $x_1, x_2$ and the corresponding $x_1', x_2'$, we obtain
\[
(x_1x_2) = \frac{r\,(x_1'x_2')}{(\gamma x_1'+\delta)(\gamma x_2'+\delta)},
\]
and if we take in addition a second pair of points $x_3, x_4$ and the corresponding $x_3', x_4'$:
\[
(x_1x_2)(x_3x_4) = \frac{r^2(x_1'x_2')(x_3'x_4')}{(\gamma x_1'+\delta)(\gamma x_2'+\delta)(\gamma x_3'+\delta)(\gamma x_4'+\delta)}.
\]
If we here interchange $x_2$ with $x_3$ and form the quotient, there follows
\weq{3}{\frac{(x_1x_2)(x_3x_4)}{(x_1x_3)(x_2x_4)} = \frac{(x_1'x_2')(x_3'x_4')}{(x_1'x_3')(x_2'x_4')}.}

The expression on the left side is called the cross ratio of the pair of points $x_1x_4$ to the pair of points $x_2x_3$, and this formula therefore contains the theorem:

\gesperrt{The cross ratio of two pairs of points is invariant under simultaneous linear transformation of the four points.}

Four points can be decomposed into two pairs in six ways. But if we set
\begin{align*}
a &= (x_2x_3)(x_1x_4) \\
b &= (x_3x_1)(x_2x_4), \\
c &= (x_1x_2)(x_3x_4)
\end{align*}
%CONT>>
